\chapter{Abdomen}

\section*{Patient focus}
This chapter explains the examination or clinical problem in accessible language. Patients should use it to prepare questions and understand the purpose, experience, benefits, and limitations of imaging; local instructions and the official report take priority.


\label{ch:ch44_abdomen}

When you have abdominal pain, your doctor may order an imaging test to find the cause. The most common investigations are \textbf{ultrasound} (which uses sound waves and is often the first test for gallstones, kidney stones, or appendicitis in children), \textbf{CT of the abdomen} (which gives detailed cross-sectional images of all the organs), and occasionally \textbf{MRI} (which is especially useful for the bile ducts and certain liver conditions).

For a CT of the abdomen, you may be asked to drink contrast dye (oral contrast) to highlight the bowel, and you will receive an injection of contrast through a vein (intravenous contrast) to assess the organs and blood vessels. The scan takes only a few minutes. Oral contrast is typically given 1--2 hours before the scan to allow it to reach the bowel. If you are having the scan in an emergency, oral contrast may be omitted to save time.

CT is the most important imaging test for acute abdominal conditions---it can detect appendicitis, diverticulitis, bowel obstruction, free air (which indicates a perforated bowel), bleeding, and many other problems. Your radiologist will examine every organ systematically.

---

\section{Coverage}
Abdominal imaging includes radiography, ultrasound, CT, MRI, and nuclear studies. Protocols should be tailored to the suspected organ, phase of enhancement, need for enteric contrast, and patient stability.

\begin{table}[h]
\centering
\caption{CT protocol for acute abdomen}
\begin{tabular}{p{3cm} p{4.5cm} p{4.5cm}}
\hline
\textbf{Component} & \textbf{Details} & \textbf{Purpose} \\
\hline
Delayed phase & 5--10 min post-injection & Urinary tract (CT urography), bladder tumours \\
IV contrast (arterial phase) & Scan at 25--35 s post-injection & Detect active haemorrhage (extravasation), vascular pathology \\
IV contrast (portal venous phase) & Scan at 60--70 s post-injection & Standard phase for solid organ assessment \\
Non-contrast & Pre-contrast images & Detect calcification (stones), acute haemorrhage, fat \\
Oral contrast & Water-soluble contrast (Gastrografin) or positive contrast (dilute barium) given 1--2 h before scan & Opacify bowel lumen; detect leaks and obstruction \\
\hline
\end{tabular}
\end{table}

\begin{pearlbox}[Oral contrast in emergencies]
In the acute setting, oral contrast is often omitted because it delays scanning by 1--2 hours. IV contrast alone is sufficient for most acute abdominal diagnoses. Oral contrast is reserved for suspected bowel leak, fistula, or pre-operative planning.
\end{pearlbox}

---

\section{Systematic Review}
Review liver, biliary tree, spleen, pancreas, adrenals, kidneys, bowel, mesentery, peritoneum, retroperitoneum, vessels, abdominal wall, and bones. Use prior studies and describe both the primary finding and complications.

---

\section{Tracing Bowel on CT}
A systematic approach to bowel assessment on CT:

\begin{enumerate}
  \item \textbf{Identify the bowel}: Start at the gastro-oesophageal junction and trace the stomach $\rightarrow$ duodenum $\rightarrow$ jejunum (left upper quadrant, valvulae conniventes) $\rightarrow$ ileum (right lower quadrant, smooth wall) $\rightarrow$ caecum $\rightarrow$ ascending $\rightarrow$ transverse $\rightarrow$ descending $\rightarrow$ sigmoid $\rightarrow$ rectum.
  \item \textbf{Wall thickness}: Normal small bowel wall $< 3$ mm (distended); normal colon $< 5$ mm. Wall thickening ($> 3$ mm small bowel, $> 5$ mm colon) indicates inflammation, ischaemia, infection, or tumour.
  \item \textbf{Enhancement}: Mural hyperenhancement suggests active inflammation (e.g., appendicitis, diverticulitis). Mural hypo-enhancement suggests ischaemia.
  \item \textbf{Calibre}: Small bowel $> 3$ cm is dilated; colon $> 6$ cm is dilated (caecum $> 9$ cm risks perforation).
  \item \textbf{Fat stranding}: Hazy increased attenuation in the mesenteric fat adjacent to abnormal bowel indicates active inflammation.
\end{enumerate}

---

\section{Bowel Obstruction}
Bowel obstruction is one of the most common acute surgical presentations.

\begin{table}[h]
\centering
\caption{Bowel obstruction on CT}
\begin{tabular}{p{3cm} p{4.5cm} p{4.5cm}}
\hline
\textbf{Finding} & \textbf{Small bowel obstruction (SBO)} & \textbf{Large bowel obstruction (LBO)} \\
\hline
Closed-loop sign & U-shaped dilated loop with twisted mesentery & Urgent (risk of strangulation) \\
Dilated loops & Small bowel $> 3$ cm with air--fluid levels & Colon $> 6$ cm; caecum $> 9$ cm \\
Pneumatosis intestinalis & Gas in bowel wall $\rightarrow$ ischaemia & Urgent \\
Small bowel faeces sign & Particulate matter in dilated small bowel & N/A \\
Transition point & Sudden change from dilated to collapsed bowel & Look for obstructing mass (sigmoid, caecal) \\
Venous pneumatosis / portal venous gas & Gas in mesenteric/portal veins $\rightarrow$ advanced ischaemia & Emergent \\
\hline
\end{tabular}
\end{table}

\begin{pearlbox}[Common causes of SBO]
Post-surgical adhesions (60\%) are the most common cause of SBO. Hernias (15\%) are the second most common and can cause closed-loop obstruction. Neoplasms, intussusception, and volvulus account for the remainder.
\end{pearlbox}

---

\section{Free Air (Pneumoperitoneum)}
Free intraperitoneal air indicates bowel perforation until proven otherwise.

\begin{itemize}
  \item \textbf{Upright CXR}: Air under the diaphragm (most sensitive plain film view).
  \item \textbf{CT abdomen}: Small bubbles of air anterior to the liver (supine patient); Rigler's sign (air on both sides of the bowel wall).
  \item \textbf{Causes}: Perforated peptic ulcer (most common), perforated diverticulitis, perforated appendicitis, post-surgical, ischaemic bowel perforation.
  \item \textbf{Surgical emergency}: Requires urgent surgical consultation.
\end{itemize}

---

\section{Appendicitis}
In nonpregnant adults, CT is highly accurate for appendicitis and is often used when imaging is indicated; ultrasound or MRI may be preferred in some patients and settings. Accuracy varies with protocol, patient factors and interpretation.

Key CT findings:
\begin{itemize}
  \item Enlarged appendix, commonly greater than 6 mm, interpreted together with wall changes and periappendiceal findings; diameter alone is not diagnostic.
  \item \textbf{Appendiceal wall enhancement} (hyperenhancing)
  \item \textbf{Appendicolith} (calcified faecolith in 25--30\%)
  \item \textbf{Peri-appendiceal fat stranding}
  \item \textbf{Target sign} (axial view): Concentric rings of the inflamed appendix
  \item \textbf{Complications}: Perforation (abscess, free fluid), phlegmon, appendicolith extrusion
\end{itemize}

---

\section{Diverticulitis}
CT findings of acute diverticulitis:
\begin{itemize}
  \item \textbf{Colonic wall thickening} ($> 4$ mm) with adjacent \textbf{fat stranding}---the hallmark finding
  \item \textbf{Diverticula}: Outpouchings of the colonic wall (most commonly sigmoid)
  \item \textbf{Abscess}: Rim-enhancing fluid collection (Hinchey classification)
  \item \textbf{Perforation}: Localised or free extraluminal air
  \item \textbf{Fistula}: Colovesical fistula (air in the bladder), colovaginal fistula
\end{itemize}

\begin{pearlbox}[Diverticulitis vs colon cancer]
Focal wall thickening with fat stranding and visible diverticula is classic for diverticulitis. However, if wall thickening is asymmetric, mass-like, or lymphadenopathy is present, consider underlying malignancy. Recommend follow-up colonoscopy 6--8 weeks after resolution of acute diverticulitis.
\end{pearlbox}

---

\section{Abdominal Masses and Incidental Findings}
When an incidental abdominal mass is detected, the radiologist should describe:
\begin{itemize}
  \item \textbf{Location} and organ of origin
  \item \textbf{Size} (longest diameter in cm)
  \item \textbf{Attenuation/signal}: Solid, cystic, mixed, fat-containing, calcified
  \item \textbf{Enhancement pattern}: Homogeneous vs heterogeneous; degree relative to cortex
  \item \textbf{Margin}: Well-defined vs ill-defined; invasion of adjacent structures
  \item \textbf{Lymphadenopathy}: Regional and distant
  \item \textbf{Staging} if malignancy suspected
\end{itemize}

---

\section{Hernias}
CT reliably detects and characterises abdominal wall hernias:
\begin{itemize}
  \item \textbf{Inguinal hernia}: Most common; indirect (through deep inguinal ring) vs direct (through Hesselbach's triangle)
  \item \textbf{Femoral hernia}: Through the femoral canal; higher risk of strangulation; more common in women
  \item \textbf{Incisional hernia}: At prior surgical sites
  \item \textbf{Complex hernias}: Spigelian, lumbar (Petit's triangle), obturator
  \item \textbf{Incarceration}: Herniated contents that cannot be reduced; bowel obstruction
  \item \textbf{Strangulation}: Ischaemic bowel within a hernia (wall thickening, hypo-enhancement, fat stranding)---surgical emergency
\end{itemize}

---

\section{Structured Report: Acute Abdomen}

\begin{testbox}[Indications:]
45-year-old female with right iliac fossa pain, nausea, and low-grade fever (37.8$^\circ$C) for 18 hours. WBC 13.2.
\end{testbox}
\begin{refbox}[Findings:]
\textbf{CT abdomen and pelvis with IV contrast (portal venous phase)}:
\textbf{Appendix}: The appendix is visualised at the base of the caecum. It measures 11 mm in diameter (normal $< 6$ mm) with circumferential wall hyperenhancement and peri-appendiceal fat stranding. A 4 mm appendicolith is identified within the lumen. No perforation, abscess, or free fluid is seen.\\
\textbf{Bowel}: Remaining bowel is unremarkable. No obstruction. No colonic wall thickening.\\
\textbf{Solid organs}: Liver, spleen, pancreas, and kidneys are normal. No hydronephrosis.\\
\textbf{Pelvis}: No free fluid. Urinary bladder is unremarkable.\\
\textbf{Bones}: No lytic or sclerotic lesions.
\end{refbox}
\begin{pearlbox}[Impression:]
1. Acute appendicitis with appendicolith. No perforation or abscess formation.\\
2. Recommend surgical consultation for appendicectomy.
\end{pearlbox}

---

\section{Safety}
Pregnancy, renal function, contrast, fasting, and sedation are protocol-dependent. Free intraperitoneal air, active bleeding, ischemia, obstruction with compromise, and ruptured aneurysm require urgent communication.

\clearpage
